% Options for packages loaded elsewhere
\PassOptionsToPackage{unicode}{hyperref}
\PassOptionsToPackage{hyphens}{url}
\PassOptionsToPackage{dvipsnames,svgnames,x11names}{xcolor}
\documentclass[
  11pt,
]{article}
\usepackage{xcolor}
\usepackage[margin=1in]{geometry}
\usepackage{amsmath,amssymb}
\setcounter{secnumdepth}{-\maxdimen} % remove section numbering
\usepackage{iftex}
\ifPDFTeX
  \usepackage[T1]{fontenc}
  \usepackage[utf8]{inputenc}
  \usepackage{textcomp} % provide euro and other symbols
\else % if luatex or xetex
  \usepackage{unicode-math} % this also loads fontspec
  \defaultfontfeatures{Scale=MatchLowercase}
  \defaultfontfeatures[\rmfamily]{Ligatures=TeX,Scale=1}
\fi
\usepackage{lmodern}
\ifPDFTeX\else
  % xetex/luatex font selection
\fi
% Use upquote if available, for straight quotes in verbatim environments
\IfFileExists{upquote.sty}{\usepackage{upquote}}{}
\IfFileExists{microtype.sty}{% use microtype if available
  \usepackage[]{microtype}
  \UseMicrotypeSet[protrusion]{basicmath} % disable protrusion for tt fonts
}{}
\makeatletter
\@ifundefined{KOMAClassName}{% if non-KOMA class
  \IfFileExists{parskip.sty}{%
    \usepackage{parskip}
  }{% else
    \setlength{\parindent}{0pt}
    \setlength{\parskip}{6pt plus 2pt minus 1pt}}
}{% if KOMA class
  \KOMAoptions{parskip=half}}
\makeatother
\usepackage{color}
\usepackage{fancyvrb}
\newcommand{\VerbBar}{|}
\newcommand{\VERB}{\Verb[commandchars=\\\{\}]}
\DefineVerbatimEnvironment{Highlighting}{Verbatim}{commandchars=\\\{\}}
% Add ',fontsize=\small' for more characters per line
\newenvironment{Shaded}{}{}
\newcommand{\AlertTok}[1]{\textcolor[rgb]{1.00,0.00,0.00}{\textbf{#1}}}
\newcommand{\AnnotationTok}[1]{\textcolor[rgb]{0.38,0.63,0.69}{\textbf{\textit{#1}}}}
\newcommand{\AttributeTok}[1]{\textcolor[rgb]{0.49,0.56,0.16}{#1}}
\newcommand{\BaseNTok}[1]{\textcolor[rgb]{0.25,0.63,0.44}{#1}}
\newcommand{\BuiltInTok}[1]{\textcolor[rgb]{0.00,0.50,0.00}{#1}}
\newcommand{\CharTok}[1]{\textcolor[rgb]{0.25,0.44,0.63}{#1}}
\newcommand{\CommentTok}[1]{\textcolor[rgb]{0.38,0.63,0.69}{\textit{#1}}}
\newcommand{\CommentVarTok}[1]{\textcolor[rgb]{0.38,0.63,0.69}{\textbf{\textit{#1}}}}
\newcommand{\ConstantTok}[1]{\textcolor[rgb]{0.53,0.00,0.00}{#1}}
\newcommand{\ControlFlowTok}[1]{\textcolor[rgb]{0.00,0.44,0.13}{\textbf{#1}}}
\newcommand{\DataTypeTok}[1]{\textcolor[rgb]{0.56,0.13,0.00}{#1}}
\newcommand{\DecValTok}[1]{\textcolor[rgb]{0.25,0.63,0.44}{#1}}
\newcommand{\DocumentationTok}[1]{\textcolor[rgb]{0.73,0.13,0.13}{\textit{#1}}}
\newcommand{\ErrorTok}[1]{\textcolor[rgb]{1.00,0.00,0.00}{\textbf{#1}}}
\newcommand{\ExtensionTok}[1]{#1}
\newcommand{\FloatTok}[1]{\textcolor[rgb]{0.25,0.63,0.44}{#1}}
\newcommand{\FunctionTok}[1]{\textcolor[rgb]{0.02,0.16,0.49}{#1}}
\newcommand{\ImportTok}[1]{\textcolor[rgb]{0.00,0.50,0.00}{\textbf{#1}}}
\newcommand{\InformationTok}[1]{\textcolor[rgb]{0.38,0.63,0.69}{\textbf{\textit{#1}}}}
\newcommand{\KeywordTok}[1]{\textcolor[rgb]{0.00,0.44,0.13}{\textbf{#1}}}
\newcommand{\NormalTok}[1]{#1}
\newcommand{\OperatorTok}[1]{\textcolor[rgb]{0.40,0.40,0.40}{#1}}
\newcommand{\OtherTok}[1]{\textcolor[rgb]{0.00,0.44,0.13}{#1}}
\newcommand{\PreprocessorTok}[1]{\textcolor[rgb]{0.74,0.48,0.00}{#1}}
\newcommand{\RegionMarkerTok}[1]{#1}
\newcommand{\SpecialCharTok}[1]{\textcolor[rgb]{0.25,0.44,0.63}{#1}}
\newcommand{\SpecialStringTok}[1]{\textcolor[rgb]{0.73,0.40,0.53}{#1}}
\newcommand{\StringTok}[1]{\textcolor[rgb]{0.25,0.44,0.63}{#1}}
\newcommand{\VariableTok}[1]{\textcolor[rgb]{0.10,0.09,0.49}{#1}}
\newcommand{\VerbatimStringTok}[1]{\textcolor[rgb]{0.25,0.44,0.63}{#1}}
\newcommand{\WarningTok}[1]{\textcolor[rgb]{0.38,0.63,0.69}{\textbf{\textit{#1}}}}
\usepackage{longtable,booktabs,array}
\newcounter{none} % for unnumbered tables
\usepackage{calc} % for calculating minipage widths
% Correct order of tables after \paragraph or \subparagraph
\usepackage{etoolbox}
\makeatletter
\patchcmd\longtable{\par}{\if@noskipsec\mbox{}\fi\par}{}{}
\makeatother
% Allow footnotes in longtable head/foot
\IfFileExists{footnotehyper.sty}{\usepackage{footnotehyper}}{\usepackage{footnote}}
\makesavenoteenv{longtable}
\setlength{\emergencystretch}{3em} % prevent overfull lines
\providecommand{\tightlist}{%
  \setlength{\itemsep}{0pt}\setlength{\parskip}{0pt}}
\usepackage{xurl,needspace,etoolbox}
\setlength{\emergencystretch}{3em}
\AtBeginEnvironment{longtable}{\small}
\pretocmd{\section}{\Needspace{7\baselineskip}}{}{}
\widowpenalty=10000
\clubpenalty=10000
\allowdisplaybreaks
\usepackage{bookmark}
\IfFileExists{xurl.sty}{\usepackage{xurl}}{} % add URL line breaks if available
\urlstyle{same}
\hypersetup{
  pdftitle={An Improved Exponent Bound for the Exact Discrete Fourier Transform},
  colorlinks=true,
  linkcolor={blue},
  filecolor={Maroon},
  citecolor={Blue},
  urlcolor={blue},
  pdfcreator={LaTeX via pandoc}}

\title{An Improved Exponent Bound for the Exact Discrete Fourier
Transform}
\usepackage{etoolbox}
\makeatletter
\providecommand{\subtitle}[1]{% add subtitle to \maketitle
  \apptocmd{\@title}{\par {\large #1 \par}}{}{}
}
\makeatother
\subtitle{Research draft: round-six complex-network transfer}
\author{}
\date{October 8, 2026}

\begin{document}
\maketitle

Under the exact-complex-arithmetic, specified-root, and logarithmic-word
address model of OpenAI's \emph{An explicit power saving for the exact
discrete Fourier transform} {[}O{]}, the proposed bound is

\[
\boxed{T(n)=O\!\left(n(\log n)^{1-\delta}\right),\qquad
\delta=\frac{73}{10^6}=7.3\times10^{-5}.}
\]

This draft transfers the attributed round-six \textbf{complex}
construction of Swapnil Jain {[}J{]}, including prior whole-residual
batching {[}B{]}, two-stage topology {[}P{]}, and copied-centre
scheduling {[}C{]}, into {[}O{]}'s uniform array model. It does not
infer a Fourier bound from an integer-multiplication running time. The
finite network is pinned to commit
\texttt{f2176bc1124821bf17eb63725bd366d7bdc020a3} of
\texttt{Swapnil-jain/integer-mult-kappa}.

\textbf{Status.} This is a proposed research result supported by a
written transfer argument and reproducible exact finite checks. Neither
this Fourier theorem nor the complete external construction has been
independently reviewed or formally verified here. The external Lean file
checks numerical certificates and their integer-multiplication assembly;
it does not formalize the Fourier transfer below. No novelty or priority
claim is made for the imported constructions.

The earlier three-stage result, with proposed saving \texttt{5.5e-10},
is preserved unchanged in \texttt{archive/three-stage/}. The separate
short proof retains the simpler \texttt{5.2e-10} and \texttt{5.3e-10}
fallbacks. Their proofs do not depend on this update.

\subsection{1. The tensor bound sufficient for the Fourier
transfer}\label{the-tensor-bound-sufficient-for-the-fourier-transfer}

Use the same fixed kernel as {[}O{]}:

\[
 C=\frac12\begin{pmatrix}1+i&1-i\\1-i&1+i\end{pmatrix},
 \qquad C^2=X=\begin{pmatrix}0&1\\1&0\end{pmatrix},
 \qquad C^{-1}=XC.
\]

The target intermediate result is an algorithm applying \(C\) on all
\(k\) binary axes of an arbitrary complex array, in

\[
 O\!\left(2^k(k+1)^\theta\right),\qquad
 \theta=1-a,\qquad a=\frac{36926111}{500000000000}.
\]

The cost must include scalar arithmetic, finite-table preparation, index
preparation, array initialization, and data movement. The algorithm must
work for every integer \(k\), and with additional spectator axes.
Section 8 proves this statement from the finite network and its strict
moment certificate. Section 9 then substitutes it into {[}O{]},
Proposition 4.2 and Section 5.4.

All finite network choices below are fixed independently of the
requested Fourier length. Only rational constants and \texttt{i} are
used by the new tensor routine. On variable data its operations are
linear: addition, subtraction, and multiplication by prepared constants.
The supplied root of order less than \texttt{1024\ n\^{}3} from {[}O{]}
remains sufficient. No bit-complexity, coefficient-size, or
numerical-stability conclusion is asserted.

\subsection{2. The pinned scalar
producer}\label{the-pinned-scalar-producer}

Fix \(h=24\) and put

\[
 \mathcal T=\binom{[h]}3,\quad v=2024,\quad m=h^2=576,\quad N=v^2=4096576.
\]

{[}J{]}'s \texttt{NStar3} producer is a finite addition DAG on inputs
\(x_T\), indexed by triples. It shares pair-exclusion sums and emits
signed partial sums whose aggregate is

\[
 D_S-H_{2,S},\qquad
 D_S=\sum_{T\cap S=\varnothing}x_T,\quad
 H_{2,S}=\sum_{|T\cap S|=2}x_T.
\]

It also retains the \(h\) totals

\[
 E_j=\sum_{T:j\notin T}x_T\quad(0\le j<h-1),
 \qquad A=\sum_Tx_T.
\]

The omitted total is recovered by

\[
 E_{h-1}=(h-3)A-\sum_{j<h-1}E_j.
\]

For a target triple \(S\), define the retained-total readout

\[
 R_S=A-\frac12\sum_{j\in S}E_j.
\]

Its coefficient on \(x_T\) is \((|S\cap T|-1)/2\). Therefore

\[
 R_S+\frac12D_S-\frac12H_{2,S}=x_S.
 \tag{2.1}
\]

Indeed the four possible intersection sizes \(0,1,2,3\) give
coefficients \(0,0,0,1\). The expression for the omitted \(E\) is used
only in this source-dependent identity. No relation among arbitrary old
scratch values is assumed.

For reproducibility, the finite object is specified by the
\textbf{unmodified, hash-pinned} producer in
\texttt{verification/vendor/jain\_round6/producer.py}. The new checker
rebuilds every support from the DAG's arguments and checks all
\(v^2=4{,}096{,}576\) entries of (2.1), including zero entries. It also
checks every retained total. This is exhaustive finite arithmetic at the
selected size, not a proof for every possible ground size.

The active producer has \texttt{28,944} addition nodes, \texttt{20,240}
signed output uses, and \texttt{24} retained output uses. Its role
compiler allocates one carrier per directed use and identifies one input
carrier with the pivot output of each addition. Thus the number of
physical auxiliary roles in one invocation is

\[
 R=28944+20240+24=49208.
 \tag{2.2}
\]

Every addition and fan-out is implemented by reversible additions into
arbitrary existing values. Let \(L\) denote this invertible scratch
mixer, \(V\) the source injection, and \(J\) the combined half-weighted
piece and retained-total readout. Equation (2.1) is \(JLV=I\). Starting
with arbitrary scratch \(z\), the chronological word

\[
 L,\ -J,\ L^{-1},\ +V,\ L,\ +J,\ L^{-1},\ -V
 \tag{2.3}
\]

changes the target by \(JL(z+Vx)-JLz=x\) and restores \(z\). Here the
signs on \(J\) and \(V\) specify additions to the target and scratch,
respectively. Reversing the word and each gate implements the inverse
shear. The scalar correctness of this construction holds on every
scratch value.

\subsection{3. Phase frames and the two
stages}\label{phase-frames-and-the-two-stages}

All label geometry is over the binary field; payloads remain complex.
For a nondegenerate binary subspace \(U\) of the ambient
\(m\)-dimensional space, use {[}O{]}'s frame

\[
 \Phi_U=H_m\operatorname{diag}_y\!\left(i^{\operatorname{wt}(P_Uy)}\right)H_m,
\]

where \(P_U\) is orthogonal projection and \(H_m\) is the normalized
Walsh matrix. The Walsh matrices establish identities and are
\textbf{not executed} by the algorithm. A scalar gate commutes with a
common address frame on its touched roles. Consequently the inserted
edge quotients telescope exactly as in {[}O{]}, equation (2.6).

The local binary space is \(D=\mathbb F_2^h\). Each source triple is
labelled by its indicator line. For a producer node, use the span of its
triple indicators when they share a common pair; otherwise use the
coordinate space of their union. The common-pair vectors are
orthonormal. These are {[}J{]}'s labels, with its output-piece splitting
retained.

The following forward schedule implements (2.3). \(J_{\rm ret}\) and
\(J_{\rm piece}\) are the two parts of \(J\); a source or target line is
the line of that particular triple.

\Needspace{24\baselineskip}

{\def\LTcaptype{none} % do not increment counter
\begin{longtable}[]{@{}
  >{\raggedright\arraybackslash}p{(\linewidth - 4\tabcolsep) * \real{0.3333}}
  >{\raggedright\arraybackslash}p{(\linewidth - 4\tabcolsep) * \real{0.3333}}
  >{\raggedright\arraybackslash}p{(\linewidth - 4\tabcolsep) * \real{0.3333}}@{}}
\toprule\noalign{}
\begin{minipage}[b]{\linewidth}\raggedright
Step
\end{minipage} & \begin{minipage}[b]{\linewidth}\raggedright
Scalar operation
\end{minipage} & \begin{minipage}[b]{\linewidth}\raggedright
Local common label
\end{minipage} \\
\midrule\noalign{}
\endhead
\bottomrule\noalign{}
\endlastfoot
1 & Apply \(L\) & \(0\) \\
2 & Subtract retained readout from target & \(0\) \\
3 & Subtract piece readout from target & \(0\) \\
4 & Apply \(L^{-1}\) & \(0\) \\
5 & Add sources to their input carriers & Source line \\
6 & Apply \(L\) & Producer-node label at each gate \\
7 & Add retained readout to target & \(0\), supplied by copies as in
Section 4 \\
8 & Add piece readout to target & Target-line perpendicular space \\
9 & Apply \(L^{-1}\) & \(D\) \\
10 & Subtract sources from their input carriers & \(D\) \\
\end{longtable}
}

The inverse schedule reverses the operations, inverts their scalar
gates, and complements every local label in \(D\). It is not a reversal
of the forward labels without complementation.

For a data pair \((S,T)\), write \(P\) and \(Q\) for its two indicator
lines. Stage one performs the forward shear on the first triple index,
for every fixed second index, with lifted label \(U\otimes Q\). Stage
two performs the inverse shear on the second triple index, with logical
source and target exchanged, and with lifted local label

\[
 (P^\perp\otimes D)\ \perp\ (P\otimes U^\perp),
 \tag{3.1}
\]

where \(U\) is the corresponding forward label before complementation.
There are \(v\) invocations per stage, with \textbf{distinct} auxiliary
banks for the two stages.

The input data frames are \(\Phi_{P\otimes Q}\) on \(X\) and identity on
\(Y\). Their terminal frames are \(\Phi_{D\otimes D}\) on \(X\) and
\(\Phi_{(P\otimes Q)^\perp}\) on \(Y\). Each auxiliary has source frame
identity and sink frame \(\Phi_{D\otimes D}\). Stage-one auxiliaries pay
an exit of rank \(m-h\); stage-two auxiliaries pay an entrance of rank
\(m-h\). No translated auxiliary-source gauge is needed for this
particular histogram.

The two interstage data residuals are both \(P^\perp\otimes Q^\perp\),
of rank \((h-1)^2\). Each of the four remaining data-front edges has
rank \(h-1\). A coordinate unit outside each triple supplies a norm-one
vector in the corresponding complement; tensoring the bases preserves
nondegeneracy and nonalternation.

For the local circuit, the verifier constructs symmetric binary
projectors and replays every physical auxiliary edge in both
orientations. It checks nesting and the existence of an orthonormal
residual basis. It separately reproduces {[}J{]}'s \texttt{140,064}
label checks with zero bad edges. Thus the new construction stays within
{[}O{]}, Lemma 2.2; it does not require the more general
alternating-residual normal form used by some other follow-ups.

\subsection{4. Copied retained centres, with their transformations
charged}\label{copied-retained-centres-with-their-transformations-charged}

The retained slots are distinct terminal output-use carriers. After the
middle producer, they are only read by the scatter before cleanup. Let
\(U\) be one such carrier's local label and \(r=\dim U\).

In the forward invocation, its old path was

\[
 U\longrightarrow0\longrightarrow D,
\]

at cost \(r+h\). Copy the physical value at frame \(U\), transform the
copy to frame zero, supply all scatter reads from the copy, and discard
it. The original remains at \(U\) until its direct move to \(D\), of
rank \(h-r\). The copied transform is an inverse residual of rank \(r\).
The new cost is \(r+(h-r)=h\), a saving of \(r\).

In the reverse complementary invocation the old path is
\(0\to D\to U^\perp\). Advance the original from zero to \(U^\perp\), at
rank \(h-r\); copy it and transform the copy to \(D\), at rank \(r\),
for the early readout. The original already has the frame required by
the subsequent inverse mixer. In both orientations, the copy's
transformed value is exactly the value the old readout would have seen.
Its scalar reads do not modify it, and all later original incidences are
preserved.

This is {[}C{]}'s copied-centre argument, specialized to {[}J{]}'s
retained totals. There are \(h-1\) retained labels of dimension \(h-1\)
and one of dimension \(h\). The loss per local invocation is therefore

\[
 \ell=(h-1)^2+h=553.
\]

After the change, the internal residual histogram \(H_r\) has mass

\[
 \sum_r rH_r=hR+\ell.
 \tag{4.1}
\]

Every copy transform remains in \(H_r\). The original and copied values
occupy full streams with all spectator coordinates. A fresh temporary
need not be an additional arbitrary-input role: copying, reading and
discarding it cost linear work, and its recursive transform is
explicitly counted. One temporary can be reused by processing retained
readouts sequentially. This is workspace, not an uncharged enlargement
of the role count \(W\).

\subsection{5. The paid endpoint
correction}\label{the-paid-endpoint-correction}

Two shears give the scalar data map

\[
 (x,y)\longmapsto(-y,x+y).
\]

For one data pair let \(u\) generate \(P\otimes Q\). It has binary norm
one and integer weight nine. Set

\[
 T=\Phi_{\langle u\rangle},\qquad F=\Phi_{D\otimes D},\qquad
 E=\Phi_{\langle u\rangle^\perp}=FT^{-1}.
\]

Starting with physical inputs \(x,y\), the two-stage common-frame
identity gives physical outputs

\[
 A=-Fy,\qquad B=FT^{-2}x+Ey.
 \tag{5.1}
\]

Apply \(T^{-1}\) to a \textbf{copy} of \(A\), and add that result to
\(B\). Since \(T^{-1}A=-Ey\), the corrected \(B\) is \(FT^{-2}x\). This
is one rank-one child per data pair, with copying, addition and erasure
charged as linear work.

A final address translation suffices to normalize this endpoint in the
array model. From the definition of the line frame,

\[
 T^2=R_u,\qquad R_u^2=I,
\]

where \(R_u\) translates binary addresses by \(u\). All these operators
commute. Thus \(R_uB=Fx\) after the correction, while \(-A=Fy\).
Reordering the banks gives the desired \(F\) on both original data
arrays. This is the same translation identity used in {[}O{]},
Proposition 2.4, applied to the weight-nine line of the two-stage
construction.

Every auxiliary scalar value is restored by its invocation. Its
source/sink frame quotient is \(F\), so it too receives exactly \(F\),
on arbitrary input data. The argument tensors over every column \(f\)
and is unchanged by spectator bits. The correction has rank \(f\) after
tensor lifting and is charged as a rank-one child on \(f\) active axes,
once for each of the \(N\) data pairs.

\subsection{6. Whole-residual recursive calls and the physical
histogram}\label{whole-residual-recursive-calls-and-the-physical-histogram}

For a nested edge with orthonormal residual basis \(z_1,\ldots,z_r\),
{[}O{]}, Lemma 2.2 writes its operator as the product of \(r\) forward
or inverse directional kernels. Use \textbf{one} binary basis extension
placing these \(r\) independent directions in the first \(r\) coordinate
slots. On \(f\) columns, collect their \(rf\) bits into one contiguous
field.

In these coordinates all factors act on disjoint bits. The identity
\(C^{-1}=XC\) moves every inverse into a translation of its bit, so the
entire residual is one \(C^{\otimes rf}\) call with address permutations
and translations before or after it. In particular, the directions'
signs need not agree. This is whole-residual batching {[}B{]}
specialized to the array primitives already allowed by {[}O{]}. The cost
of the wrappers is established in Section 8, rather than assumed free.

The full list of child widths consists of the following disjoint
classes:

\Needspace{12\baselineskip}

{\def\LTcaptype{none} % do not increment counter
\begin{longtable}[]{@{}lrr@{}}
\toprule\noalign{}
Class & Number of children & Width \\
\midrule\noalign{}
\endhead
\bottomrule\noalign{}
\endlastfoot
Internal auxiliary edges, including centre copies & \(2vH_r\) & \(r\) \\
Auxiliary exterior edges & \(2vR\) & \(m-h=552\) \\
Interstage data edges & \(2N\) & \((h-1)^2=529\) \\
Data fronts & \(4N\) & \(h-1=23\) \\
Endpoint copies & \(N\) & \(1\) \\
\end{longtable}
}

With

\[
 W=2N+2vR=207387136,\qquad L=2v\ell=2238544,
\]

the sum of all child widths is

\[
 s=Wm-N+L=119453132304,\qquad Wm-s=1858032.
 \tag{6.1}
\]

Both orientations give the same \(H_r\) under the physical projector
replay. The resulting full histogram is:

\Needspace{20\baselineskip}

{\def\LTcaptype{none} % do not increment counter
\begin{longtable}[]{@{}rrrr@{}}
\toprule\noalign{}
\(r\) & \(n_r\) & \(r\) & \(n_r\) \\
\midrule\noalign{}
\endhead
\bottomrule\noalign{}
\endlastfoot
1 & 184811440 & 14 & 14864256 \\
2 & 107199136 & 15 & 4849504 \\
3 & 50583808 & 16 & 13014320 \\
4 & 53478128 & 17 & 10196912 \\
5 & 33031680 & 18 & 33371712 \\
6 & 27611408 & 19 & 13601280 \\
7 & 14767104 & 20 & 22774048 \\
8 & 20195472 & 22 & 44038192 \\
9 & 6217728 & 23 & 16479408 \\
10 & 14119424 & 24 & 4048 \\
11 & 5051904 & 529 & 8193152 \\
12 & 13742960 & 552 & 199193984 \\
13 & 4663296 & & \\
\end{longtable}
}

Unlisted widths have multiplicity zero. Every child has
\(1\le r\le552<576\), so there is no same-width recursion. The verifier
checks the complete table against both a fresh run of {[}J{]}'s
histogram builder and the pinned inputs to \texttt{Round6.lean}.

\subsection{7. An exact rational moment
certificate}\label{an-exact-rational-moment-certificate}

Define

\[
 \Psi(\theta)=\frac1W\sum_r n_r\left(\frac rm\right)^\theta.
\]

Use the rational saving actually checked by {[}J{]}'s Lean file,

\[
 a=\frac{36926111}{500000000000}=0.000073852222,
 \qquad \theta=1-a.
\]

For rational upper bounds \(U_r\ge\log(m/r)\),

\[
 \Psi(1-a)
 =\sum_r\frac{rn_r}{Wm}\exp\!\left(a\log\frac mr\right)
 \le\sum_r\frac{rn_r}{Wm(1-aU_r)},
 \tag{7.1}
\]

provided \(0\le aU_r<1\). The inequality follows by comparing the
exponential and geometric power series. The logarithm bounds use range
reduction by powers of two and

\[
 \log y=2\sum_{j=0}^{K-1}\frac{z^{2j+1}}{2j+1}+\mathcal R_K,
 \quad z=\frac{y-1}{y+1},\quad
 0\le\mathcal R_K\le\frac{2z^{2K+1}}{(2K+1)(1-z^2)}.
\]

The new verifier uses the repository's earlier logarithm implementation,
rounded \textbf{upward} to a rational grid, rather than importing
{[}J{]}'s logarithm routine. Exact arithmetic gives

\[
 \sum_r\frac{rn_r}{Wm(1-aU_r)}
 <1-\frac7{10^{14}}<1.
 \tag{7.2}
\]

Its actual rational gap is approximately
\(7.91950201573749\times10^{-14}\); the certificate stores the rational
endpoints, not merely this decimal. This is a strict witness, not a
claim to have attained the critical root of the moment equation.
{[}J{]}'s tighter Python bound \(7.3861113\times10^{-5}\) is not needed
here.

\subsection{8. Uniform unequal-width recursion without role
padding}\label{uniform-unequal-width-recursion-without-role-padding}

\textbf{Linear-time layout lemma.} A fixed binary basis map on each of
\(f\) columns, selection of any fixed \(r\) output slots per column, and
fixed column translations can be implemented in \(O(2^k)\) charged
operations for \(k=mf+b\), \(0\le b<m\).

To see this, extend {[}O{]}, Lemma 2.5 from one selected slot to \(r\)
slots. Treat each \(m\)-bit column as a digit of fixed radix \(2^m\).
Its finite table records contributions to the selected \(rf\)-bit field
and to the spectator fields. Precompute their displaced contributions in
\(O(f+1)\) time. Traverse the digit prefix tree, maintaining source and
destination address sums. Every extension uses a fixed number of word
operations; fewer than \(2^{k+1}\) nodes are visited. At each leaf move
the value once, using an output buffer. The inverse permutation reverses
the source/destination addresses. Translations can be incorporated into
the digit tables. This includes initialization, discovery and movement
of the contiguous fibers and requires no per-entry scan of \(f\)
columns.

Now let \(T(k)\) be the cost of applying \(C^{\otimes k}\) to \(W\)
arbitrary arrays and put \(t(k)=T(k)/(W2^k)\). Handle the fixed range
\(k<m\) directly. For \(k\ge m\), set \(f=\lfloor k/m\rfloor\) and
process the fewer than \(m\) leftover axes directly. Apply the finite
network on the remaining \(mf\) axes.

A child of width \(r\) acts on \(F_r=2^{k-rf}\) fibers of length
\(2^{rf}\) in \textbf{one} original or temporary stream. Pack complete
groups of \(W\) fibers and run this same simultaneous algorithm
recursively. Process the fewer than \(W\) leftover fibers by the
ordinary coordinate algorithm. This gives

\[
\begin{split}
 T(k)\le{}&A W2^k+\sum_r n_r\left\lfloor\frac{2^{k-rf}}W\right\rfloor T(rf)\\
 &+B(W-1)\sum_r n_r\,rf\,2^{rf}.
\end{split}
 \tag{8.1}
\]

The first term charges the fixed number of scalar gates, all layout
passes, temporary copies, endpoint translations, and leftover axes. All
constants are independent of \(k\). Since \(r<m\) and \(k\ge mf\),

\[
 rf\,2^{rf-k}\le rf\,2^{-(m-r)f}\le r.
\]

Consequently incomplete batches contribute only a bounded additive term
after normalization:

\[
 t(k)\le A'+\frac1W\sum_r n_r\,t(r\lfloor k/m\rfloor).
 \tag{8.2}
\]

Choose a constant \(K\) large enough for the finite base range and for
\(K(1-\Psi(\theta))m^\theta\ge A\prime\). Strong induction on \(k\),
using \(rf\le(r/m)k<k\), gives

\[
 t(k)\le A'+Kk^\theta\Psi(\theta)\le Kk^\theta
 \quad(k\ge m).
\]

Thus \(T(k)=O(W2^k(k+1)^\theta)\). A single requested transform
initializes \(W-1\) other arrays once and retains its own output;
because \(W\) is fixed, this adds only \(O(2^k)\) work.

The recursion can be performed serially with reusable buffers and
temporary streams. Each child array has at most half as many entries as
its parent, since \(k-rf\ge1\). The maximum axis count also contracts by
at most \(552/576\), giving \(O(\log(k+2))\) stack depth. Summing live
array sizes along the stack gives \(O(W2^k)\) workspace. Addresses,
counters and strides fit a fixed number of \(O(k+1)\)-bit words. Fixed
tables and scalar coefficients have constructive finite preparation
procedures, so their initial construction is charged as a constant.
These facts establish the uniform arbitrary-input tensor interface
required by {[}O{]}.

\subsection{9. Propagation to every Fourier
length}\label{propagation-to-every-fourier-length}

The proof of {[}O{]}, Proposition 4.2 uses its tensor routine only
through the bound \(O(2^k(k+1)^\theta)\) on arbitrary complex arrays,
including index and movement costs. Replacing that routine by Section 8
therefore gives

\[
 O\!\left(R(\ell+1)^\theta(1+\log r_{\max})^4\right)
 +\operatorname{poly}(\ell,r_{\max},\log(R+2)).
\]

The exact-width local Fourier compiler, its scalar preparation and
denominator conditions are unchanged. Its sectors still carry the same
kernel \(C\). Applying {[}O{]}'s small-prime working lengths, CRT
permutations, and charged chirp convolution, with
\(\ell=\Theta(\log n/\log\log n)\), gives

\[
 T(n)=O\!\left(n(\log n)^\theta(\log\log n)^{4-\theta}\right).
 \tag{9.1}
\]

In particular, the transform of the fixed convolution operand remains
charged. The specified root order and logarithmic-word address bounds
remain those of {[}O{]}.

Finally,

\[
 \delta=\frac{73}{10^6},\qquad
 a-\delta=\frac{426111}{500000000000}>0.
\]

For every fixed positive exponent gap, the fixed power of \(\log\log n\)
in (9.1) is eventually smaller than \((\log n)^{a-\delta}\). This proves
the proposed pure-power corollary, conditional on the specified finite
construction and upstream results. There is no integer-multiplication
bit-network constraint or extra factor-of-two loss in this transfer.

The new proposed \texttt{delta} is about \texttt{132,727} times the
earlier \texttt{5.5e-10} value and \texttt{730,000,000} times {[}O{]}'s
published \texttt{1e-13} headline. These ratios compare asymptotic
exponent savings, not running times or practical FFT performance. The
previous centre reduction and shared-sum gains are not added to the
imported construction's saving.

\subsection{10. Reproduction and proof
boundaries}\label{reproduction-and-proof-boundaries}

Run from the repository root with standard-library Python 3.10 or newer:

\begin{Shaded}
\begin{Highlighting}[]
\ExtensionTok{python3}\NormalTok{ verification/run\_checks.py}
\end{Highlighting}
\end{Shaded}

For the new component alone:

\begin{Shaded}
\begin{Highlighting}[]
\ExtensionTok{python3}\NormalTok{ verification/verify\_round6.py}
\end{Highlighting}
\end{Shaded}

The new verifier checks source hashes; reconstructs all producer
supports and the full local coefficient matrix; reproduces the external
binary label checks; replays physical auxiliary paths in both
orientations using the earlier repository's binary-projector
implementation; charges centre copies; rebuilds and compares the entire
histogram; checks dirty-scratch scalar trials and exact endpoint
identities; and certifies the rational moment and absorption slack.
Negative controls reject unpaid centre copies, an omitted endpoint
correction, and an excessive exponent saving. The report is compared
with \texttt{verification/certificates/round6\_reference.json}.

The producer and the external frame/histogram checks are vendored
source, not independent implementations. The extra replay and
coefficient checker were written with the same assistant and reuse
earlier local binary algebra. The dirty-scratch executions are finite
exact rational trials; the universal scratch statement uses (2.3) and
the copied-read proof. Endpoint tests use exact Gaussian-integer
algebra. The checker does not execute the astronomical full Fourier
algorithm, prove the infinite recurrence by enumeration, or run Lean.
The pinned Lean source and numerical inputs are included for comparison,
with their original license and notice.

Independent mathematical review should prioritize the copied-read
timing, the lifted data-frame interfaces, the paid endpoint correction,
the full-residual layout lemma, and the remainder recurrence. The
original compiler and all-length reduction are invoked with their stated
hypotheses, not re-proved here. The earlier result and its verification
suites remain separately available.

\subsection{References and
attribution}\label{references-and-attribution}

{[}O{]} OpenAI. \emph{An explicit power saving for the exact discrete
Fourier transform}, September 25, 2026. Especially Lemmas 2.2 and 2.5,
Proposition 2.4, Theorem 2.6, Proposition 4.2, and Section 5.
\href{https://github.com/openai/math/blob/main/preprints/An-explicit-power-saving-for-the-exact-discrete-Fourier-transform-September-25-2026/main.pdf}{Public
manuscript}.

{[}J{]} Swapnil Jain. \texttt{integer-mult-kappa}, round-six complex
producer and certificates, pinned at
\texttt{f2176bc1124821bf17eb63725bd366d7bdc020a3}.
\href{https://github.com/Swapnil-jain/integer-mult-kappa/tree/f2176bc1124821bf17eb63725bd366d7bdc020a3/independent/complex-twostage}{Complex
implementation},
\href{https://github.com/Swapnil-jain/integer-mult-kappa/blob/f2176bc1124821bf17eb63725bd366d7bdc020a3/lean/Round6.lean}{Lean
numerical certificate}. The source attributes research and
implementation assistance to Claude (Anthropic); its Apache-2.0 license
and full NOTICE are preserved locally.

{[}B{]} icekylinx. Whole-residual batching and unequal-width recursion,
PR \#10 of Colkitt's repository, commit \texttt{62691e3}; further full
complex batching is credited to eumemic, PR \#15.
\href{https://github.com/icekylinx/integer-mult-bounds/blob/62691e3/notes/batched-complex-rows.tex}{Batching
proof},
\href{https://github.com/CrocSwap/integer-mult-bounds/pull/15}{PR \#15}.

{[}C{]} icekylinx. Copied retained centres and two-stage complex
endpoint transfer, PR \#36, commit
\texttt{11817ccacb564bb7f98789c20dc11d3fece207e3}.
\href{https://github.com/icekylinx/integer-mult-bounds/blob/11817ccacb564bb7f98789c20dc11d3fece207e3/notes/copied-centers-lemma.tex}{Copied-centre
lemma},
\href{https://github.com/icekylinx/integer-mult-bounds/blob/11817ccacb564bb7f98789c20dc11d3fece207e3/notes/copied-centers-complex.tex}{complex
endpoint argument}. The source records substantial GPT-6 Astra and Codex
assistance.

{[}P{]} Aurel Prosz (Paureel). Two-stage topology and charged
endpoint-copy correction, commit
\texttt{c82d09eb4781b68435e3fa3eb6beea72fa900fab} of
\href{https://github.com/Paureel/integer-mult-bounds/tree/c82d09eb4781b68435e3fa3eb6beea72fa900fab}{integer-mult-bounds},
as attributed in {[}J{]} and {[}C{]}.

The common finite-network framework originates in OpenAI's \emph{Integer
multiplication below n log n} (\#109); Douglas Colkitt's
\texttt{CrocSwap/integer-mult-bounds} supplies the follow-up framework.
Other antecedents, including contributions by Zhihao Chen /
jacklightChen and retained-total and complex-circuit work, are
identified in the preserved external NOTICE. This draft claims a
proposed transfer to \#130's arithmetic model, not authorship of those
network improvements. AI-assisted checking is not independent review.

\end{document}
